%% Supplementary Information for "Reliability beyond accuracy in machine learning of magnetic materials" (Nature format, 2026-11-10).
%% Nature: Supplementary Information holds no figures (the former Fig. S1 is Extended Data Fig. 1 of the main file); tables are "Supplementary Table N".
%% Tables are local copies in tables/ (refresh with refresh_from_package.py); S9 and S10 are written here from
%% paper/results/revision/research/prior_results_table.md and data_licences.md. Compile like the main text (pdflatex, bibtex).

\documentclass[11pt,a4paper]{article}
%% Shared page design for manuscript.tex and supplementary.tex (plain article class; Nature has no required LaTeX class
%% and takes a PDF at first submission). Two modes:
%%   reading mode (default)       : single-spaced, no line numbers -- a clean copy to read and circulate;
%%   submission mode              : double-spaced with line numbers, as the Nature formatting guide asks for a PDF submission.
%% Submission mode is switched on by defining \submissionmode before this file is read (see manuscript_submission.tex).

\usepackage[a4paper,left=22mm,right=22mm,top=24mm,bottom=26mm]{geometry}
\usepackage{iftex}
\ifPDFTeX                                   % pdfLaTeX (Overleaf default)
  \usepackage[T1]{fontenc}
  \usepackage{newtxtext,newtxmath}          % Times-like text and matching mathematics
  \usepackage[scaled=0.9]{helvet}           % Helvetica-like sans serif for title, headings, captions and tables
\else                                       % XeLaTeX / LuaLaTeX / tectonic
  \usepackage{newtxtext,newtxmath}
  \setsansfont{texgyreheros}[Extension=.otf,UprightFont=*-regular,BoldFont=*-bold,ItalicFont=*-italic,BoldItalicFont=*-bolditalic,Scale=0.9]
\fi
\usepackage{graphicx}
\usepackage{xcolor}
\usepackage{amsmath}
\usepackage{booktabs,tabularx,array}
\usepackage{float}
\usepackage{setspace}
\usepackage{url}
\usepackage[super,comma,sort&compress]{natbib}   % Nature: superscript reference numbers
\usepackage{titlesec}
\usepackage{caption}
\usepackage[hidelinks]{hyperref}

\bibliographystyle{sn-nature}
\providecommand{\bibcommenthead}{}          % written by sn-nature.bst; defined by the Springer class only
\renewcommand{\bibnumfmt}[1]{#1.}           % Nature reference list: "1." not "[1]"
\setlength{\bibsep}{2pt}
\renewcommand{\bibfont}{\small}

%% Headings: bold sans serif, unnumbered, close to the text they introduce.
\titleformat{\section}{\sffamily\bfseries\large\raggedright}{}{0pt}{}
\titlespacing*{\section}{0pt}{16pt plus 2pt minus 2pt}{4pt}
\titleformat{\subsection}{\sffamily\bfseries\normalsize\raggedright}{}{0pt}{}
\titlespacing*{\subsection}{0pt}{10pt plus 2pt minus 1pt}{2pt}
\newcommand{\bmhead}[1]{\subsection*{#1}}

%% Captions in Nature style: "Fig. 1 | Title. Text" and "Table 1 | Title", small sans serif, full width.
\DeclareCaptionLabelSeparator{bar}{\ \textbar\ }
\captionsetup{font={footnotesize,sf},labelfont=bf,labelsep=bar,justification=justified,singlelinecheck=false,skip=6pt}
\captionsetup[table]{position=top,textfont=bf,skip=4pt}
\renewcommand{\figurename}{Fig.}

%% One font for every table (the table files call \tablefont).
\newcommand{\tablefont}{\sffamily\fontsize{7.5}{9.5}\selectfont}
\renewcommand{\arraystretch}{1.1}

%% Floats: allow large figures and tables at the top of a page without pushing them to the end.
\renewcommand{\topfraction}{0.9}
\renewcommand{\bottomfraction}{0.6}
\renewcommand{\textfraction}{0.08}
\renewcommand{\floatpagefraction}{0.75}
\setlength{\textfloatsep}{14pt plus 2pt minus 2pt}
\floatplacement{table}{t}
\floatplacement{figure}{t}

\setlength{\parindent}{1.2em}
\setlength{\parskip}{0pt}
\emergencystretch=2em                       % no line may stick out into the margin
\frenchspacing
\raggedbottom

\ifdefined\submissionmode
  \usepackage{lineno}
  \renewcommand{\linenumberfont}{\normalfont\scriptsize\sffamily\color{black!55}}
  \setlength{\linenumbersep}{8pt}
  \AtBeginDocument{\doublespacing\linenumbers}
\fi

\newcommand{\todo}[1]{\textcolor{red}{[TODO: #1]}}
\newcommand{\muB}{\ensuremath{\mu_\mathrm{B}}}
\newcommand{\Tc}{\ensuremath{T_\mathrm{C}}}
\newcommand{\mcell}{\ensuremath{m_\mathrm{cell}}}
\newcommand{\mfu}{\ensuremath{m_\mathrm{fu}}}

\renewcommand{\tablename}{Supplementary Table}
\floatplacement{table}{H}

\begin{document}
\pagestyle{plain}

\begin{flushleft}
{\sffamily\bfseries\fontsize{17}{21}\selectfont Supplementary Information\par}
\vspace{6pt}
{\sffamily\bfseries\large Reliability beyond accuracy in machine learning of magnetic materials\par}
\vspace{8pt}
{\sffamily Mujahidul Islam and Gayan P.~Hettiarachchi\par}
\end{flushleft}

\section*{Contents}
\begin{itemize}
\item Supplementary Tables 1--3: audit of the external Curie benchmark, subset comparisons, and error by composition distance.
\item Supplementary Table 4: the external benchmark under three resampling units and two interval types.
\item Supplementary Table 5: residual label noise in the curated Curie target.
\item Supplementary Table 6: design effects and distance-aware conformal variants.
\item Supplementary Table 7: ordering performance with and without compounds lacking a magnetism-capable element.
\item Supplementary Table 8: tier effects for the moment per formula unit.
\item Supplementary Table 9: prior ML results for magnetic ordering and Curie temperature.
\item Supplementary Table 10: data licences.
\item Supplementary Table 11: Curie-temperature interval coverage in every subgroup evaluated.
\item Supplementary Table 12: ordering labels of Materials Project against experimental sources, and transfer between them.
\item Reproducibility notes.
\end{itemize}

\begin{table}
\caption{Audit of the external Curie benchmark (2,500 records; 105 excluded in total)}
\label{tab:tableS1_benchmark_audit}
\tablefont
\begin{tabularx}{\textwidth}{@{}l>{\raggedright\arraybackslash}Xrl@{}}
\toprule
\textbf{Code} & \textbf{Check} & \textbf{Records} & \textbf{Action} \\
\midrule
A1 & Computational label (source paper reports no measurement) & 60 & excluded \\
A2 & Malformed composition (mixed-scale, impossible ionic, shorthand) & 18 & excluded \\
A3 & Duplicate records (exact or same composition) & 0 & merged \\
A4 & Outside the training target definition (rules C1-C5) & 35 & excluded \\
A5 & Transition not a magnetic Curie point (from title) & 0 & excluded \\
A6 & Possibly not directly comparable (film, nano, pressure, 2D...) & 160 & flagged \\
A7 & Composition already in training repository & 1,998 & defines subsets \\
A7 & Value identical to a value in our Curie data & 1,200 & defines subsets \\
A8 & No DOI (book reference) & 34 & flagged \\
\bottomrule
\end{tabularx}
\end{table}

\begin{table}
\caption{Grouped-bootstrap comparisons between external-benchmark subsets (first minus second subset)}
\label{tab:tableS2_subset_comparisons}
\tablefont
\begin{tabularx}{\textwidth}{@{}>{\raggedright\arraybackslash}Xlrlr@{}}
\toprule
\textbf{Clusters} & \textbf{MAE difference, K} & \textbf{Share worse (MAE)} & \textbf{R$^2$ difference} & \textbf{Share worse (R$^2$)} \\
\midrule
\multicolumn{5}{@{}l@{}}{\textit{blind+unseen system vs blind+seen system}} \\
\quad paper & +81 [$-$27, +165] & 0.89 & $-$0.95 [$-$1.69, $-$0.06] & 0.99 \\
\quad paper (first DOI) & +81 [$-$23, +164] & 0.91 & $-$0.95 [$-$1.68, $-$0.08] & 0.99 \\
\quad chemical system & +81 [+29, +145] & 1.00 & $-$0.95 [$-$1.50, $-$0.53] & 1.00 \\
\addlinespace
\multicolumn{5}{@{}l@{}}{\textit{blind vs seen composition}} \\
\quad paper & +45 [+22, +94] & 1.00 & $-$0.12 [$-$0.44, +0.01] & 0.96 \\
\quad paper (first DOI) & +45 [+23, +71] & 1.00 & $-$0.12 [$-$0.38, +0.02] & 0.95 \\
\quad chemical system & +45 [+32, +60] & 1.00 & $-$0.12 [$-$0.26, $-$0.02] & 1.00 \\
\addlinespace
\multicolumn{5}{@{}l@{}}{\textit{blind+never-used+unseen vs blind+never-used+seen}} \\
\quad paper & +39 [$-$42, +170] & 0.76 & $-$0.95 [$-$4.95, +0.06] & 0.96 \\
\quad paper (first DOI) & +39 [$-$39, +145] & 0.77 & $-$0.95 [$-$2.86, +0.05] & 0.96 \\
\quad chemical system & +39 [$-$38, +142] & 0.76 & $-$0.95 [$-$2.70, +0.05] & 0.96 \\
\bottomrule
\end{tabularx}
\end{table}

\begin{table}
\caption{Error and interval behaviour by composition distance to the training set (truly blind records)}
\label{tab:tableS3_error_vs_distance}
\tablefont
\begin{tabularx}{\textwidth}{@{}>{\raggedright\arraybackslash}Xlrlrlr@{}}
\toprule
\textbf{Distance quartile} & \textbf{Composition distance} & \textbf{n (papers)} & \textbf{MAE, K} & \textbf{Mean width, K} & \textbf{Coverage} & \textbf{Unseen-system share} \\
\midrule
Q1 (closest) & 0.000--0.015 & 117 (31) & 62.3 [49.7, 93.0] & 388 & 0.95 [0.89, 0.98] & 0.00 \\
Q2 & 0.015--0.033 & 121 (24) & 79.8 [66.4, 109.1] & 380 & 0.91 [0.86, 0.96] & 0.12 \\
Q3 & 0.034--0.083 & 114 (39) & 106.4 [81.0, 145.8] & 424 & 0.87 [0.80, 0.93] & 0.12 \\
Q4 (farthest) & 0.087--0.550 & 113 (37) & 165.3 [98.8, 240.2] & 529 & 0.85 [0.72, 0.93] & 0.36 \\
\bottomrule
\end{tabularx}
\end{table}

\begin{table}
\caption{External benchmark under three resampling units and two interval types}
\label{tab:tableS4_external_clusterings}
\tablefont
\begin{tabularx}{\textwidth}{@{}>{\raggedright\arraybackslash}Xrllll@{}}
\toprule
\textbf{Clustering} & \textbf{n (clusters)} & \textbf{R$^2$, percentile} & \textbf{R$^2$, BCa} & \textbf{Skill vs training mean, BCa} & \textbf{Upper limit $\leq$ 0.2?} \\
\midrule
\multicolumn{6}{@{}l@{}}{\textit{blind}} \\
\quad paper & 465 (81) & 0.69 [0.38, 0.81] & 0.69 [0.43, 0.82] & 0.70 [0.41, 0.84] & no \\
\quad paper (first DOI) & 465 (134) & 0.69 [0.41, 0.82] & 0.69 [0.43, 0.82] & 0.70 [0.44, 0.84] & no \\
\quad chemical system & 465 (251) & 0.69 [0.55, 0.78] & 0.69 [0.56, 0.78] & 0.70 [0.58, 0.80] & no \\
\addlinespace
\multicolumn{6}{@{}l@{}}{\textit{blind + seen chemical system}} \\
\quad paper & 396 (72) & 0.78 [0.57, 0.84] & 0.78 [0.67, 0.87] & 0.79 [0.66, 0.89] & no \\
\quad paper (first DOI) & 396 (121) & 0.78 [0.60, 0.85] & 0.78 [0.65, 0.86] & 0.79 [0.65, 0.88] & no \\
\quad chemical system & 396 (203) & 0.78 [0.68, 0.84] & 0.78 [0.70, 0.85] & 0.79 [0.71, 0.87] & no \\
\addlinespace
\multicolumn{6}{@{}l@{}}{\textit{blind + unseen chemical system}} \\
\quad paper & 69 (18) & $-$0.17 [$-$0.91, 0.68] & $-$0.17 [$-$1.02, 0.64] & $-$0.17 [$-$0.77, 0.56] & no \\
\quad paper (first DOI) & 69 (26) & $-$0.17 [$-$0.89, 0.66] & $-$0.17 [$-$0.97, 0.63] & $-$0.17 [$-$0.77, 0.56] & no \\
\quad chemical system & 69 (48) & $-$0.17 [$-$0.71, 0.23] & $-$0.17 [$-$0.75, 0.21] & $-$0.17 [$-$0.77, 0.17] & yes (R$^2$ or skill) \\
\bottomrule
\end{tabularx}
\end{table}

\begin{table}
\caption{Residual label noise in the curated Curie target}
\label{tab:tableS5_curie_label_noise}
\tablefont
\begin{tabularx}{\textwidth}{@{}>{\raggedright\arraybackslash}Xl@{}}
\toprule
\textbf{Quantity} & \textbf{Value} \\
\midrule
likely wrong label (other transition + implausible value + computational) & 22/200 = 11.0\% [7.4, 16.1] \\
non-Curie transition & 7/200 = 3.5\% [1.7, 7.0] \\
implausible value & 7/200 = 3.5\% [1.7, 7.0] \\
computational source & 8/200 = 4.0\% [2.0, 7.7] \\
paper about a different compound & 32/200 = 16.0\% [11.6, 21.7] \\
film/nano/amorphous sample & 25/200 = 12.5\% [8.6, 17.8] \\
unclear or no title & 7/200 = 3.5\% [1.7, 7.0] \\
formulas with $\geq$ 2 retained records & 3,634 (65\% of records) \\
median within-formula SD of $T_\mathrm{C}$ & 4.2 K \\
share of those formulas with SD > 50 K / > 100 K & 21\% / 13\% \\
\bottomrule
\end{tabularx}
\par\vspace{4pt}\raggedright\noindent Title-based screen, not a full-text audit: stratified random sample of 200 retained records (source x $T_\mathrm{C}$ tercile, seed 20261110), each classified by the AI assistant from the Crossref title, the compound and the plausibility of the value; Wilson 95\% intervals. Per-record labels for checking: results/revision/a7\_curie\_label\_noise/title\_audit\_adjudicated.csv.
\end{table}

\begin{table}
\caption{Why Curie interval coverage varies between splits, and distance-aware conformal variants}
\label{tab:tableS6_design_effect_distance_conformal}
\tablefont
\begin{tabularx}{\textwidth}{@{}>{\raggedright\arraybackslash}Xrrrrrr@{}}
\toprule
\textbf{Method / variant} & \textbf{Design effect} & \textbf{Effective size} & \textbf{Implied SD} & \textbf{Observed SD} & \textbf{Mean coverage} & \textbf{Farthest decile} \\
\midrule
Split conformal & 13.7 & 305 & 0.5 / 1.7 & 6.0 & -- & -- \\
Normalized conformal & 12.4 & 337 & 0.5 / 1.6 & 2.5 & -- & -- \\
CQR & 16.7 & 251 & 0.5 / 1.9 & 5.8 & -- & -- \\
Normalized, plain & -- & -- & -- & 2.5 & 90.3 & 88.2 \\
Normalized, calibrated per distance quartile & -- & -- & -- & 2.3 & 89.5 & 90.3 \\
Normalized, weighted for covariate shift & -- & -- & -- & 2.4 & 90.3 & 88.2 \\
\bottomrule
\end{tabularx}
\par\vspace{4pt}\raggedright\noindent Design effect from the within-formula intraclass correlation of the calibration exceedance indicator. Effective size = effective number of independent calibration units. Implied SD of training-conditional coverage: for exchangeable rows / for the effective size. Observed SD across 10 splits. SD and coverage in percentage points; farthest decile = coverage in the farthest decile of composition distance.
\end{table}

\begin{table}
\caption{Magnetic-ordering performance with and without compounds lacking a magnetism-capable element}
\label{tab:tableS7_ordering_decomposition}
\tablefont
\begin{tabularx}{\textwidth}{@{}>{\raggedright\arraybackslash}Xr@{}}
\toprule
\textbf{Quantity} & \textbf{Mean over 10 splits} \\
\midrule
Macro-F1 & 0.692 \\
Macro-F1 (compounds with a magnetism-capable element) & 0.681 \\
Share of test rows with a magnetism-capable element & 0.683 \\
Share of non-magnetic rows with a magnetism-capable element & 0.388 \\
Recall, FM & 0.774 \\
Recall, FM (compounds with a magnetism-capable element) & 0.800 \\
Recall, AFM & 0.546 \\
Recall, AFM (compounds with a magnetism-capable element) & 0.551 \\
Recall, FiM & 0.485 \\
Recall, FiM (compounds with a magnetism-capable element) & 0.500 \\
Recall, NM & 0.906 \\
Recall, NM (compounds with a magnetism-capable element) & 0.837 \\
Macro-F1 over FM, AFM and FiM (compounds with a magnetism-capable element) & 0.654 \\
\bottomrule
\end{tabularx}
\par\vspace{4pt}\raggedright\noindent Tier 1 XGBoost on the full ordering population, 10 prespecified splits; magnetism-capable element = curation rule C1 list.
\end{table}

\begin{table}
\caption{Tier effects for the DFT moment per formula unit and per atom}
\label{tab:tableS8_moment_per_fu}
\tablefont
\begin{tabularx}{\textwidth}{@{}>{\raggedright\arraybackslash}Xllrl@{}}
\toprule
\textbf{Algorithm} & \textbf{Contrast} & \textbf{Improvement} & \textbf{Splits better} & \textbf{Verdict} \\
\midrule
\multicolumn{5}{@{}l@{}}{\textit{per formula unit, all}} \\
\quad Ridge / logistic & Tier 1 $\rightarrow$ Tier 2 & +0.018 [+0.013, +0.023] $\mu_\mathrm{B}$/f.u. & 10/10 & stable gain \\
\quad Ridge / logistic & Tier 1 $\rightarrow$ Tier 3 & +0.053 [+0.035, +0.070] $\mu_\mathrm{B}$/f.u. & 10/10 & stable gain \\
\quad Ridge / logistic & Tier 2 $\rightarrow$ Tier 3 & +0.034 [+0.015, +0.054] $\mu_\mathrm{B}$/f.u. & 10/10 & stable gain \\
\quad Ridge / logistic & Tier 1 $\rightarrow$ Tier 2 + site count & +0.022 [+0.017, +0.027] $\mu_\mathrm{B}$/f.u. & 10/10 & stable gain \\
\quad Ridge / logistic & Tier 2 + site count $\rightarrow$ Tier 3 & +0.030 [+0.010, +0.051] $\mu_\mathrm{B}$/f.u. & 10/10 & stable gain \\
\quad Gradient boosting & Tier 1 $\rightarrow$ Tier 2 & +0.012 [$-$0.009, +0.032] $\mu_\mathrm{B}$/f.u. & 7/10 & borderline gain \\
\quad Gradient boosting & Tier 1 $\rightarrow$ Tier 3 & +0.031 [+0.014, +0.048] $\mu_\mathrm{B}$/f.u. & 10/10 & stable gain \\
\quad Gradient boosting & Tier 2 $\rightarrow$ Tier 3 & +0.019 [+0.003, +0.036] $\mu_\mathrm{B}$/f.u. & 10/10 & stable gain \\
\quad Gradient boosting & Tier 1 $\rightarrow$ Tier 2 + site count & +0.017 [$-$0.003, +0.038] $\mu_\mathrm{B}$/f.u. & 9/10 & borderline gain \\
\quad Gradient boosting & Tier 2 + site count $\rightarrow$ Tier 3 & +0.014 [$-$0.001, +0.028] $\mu_\mathrm{B}$/f.u. & 9/10 & none \\
\quad XGBoost & Tier 1 $\rightarrow$ Tier 2 & +0.009 [$-$0.013, +0.030] $\mu_\mathrm{B}$/f.u. & 7/10 & none \\
\quad XGBoost & Tier 1 $\rightarrow$ Tier 3 & +0.023 [+0.010, +0.036] $\mu_\mathrm{B}$/f.u. & 10/10 & stable gain \\
\quad XGBoost & Tier 2 $\rightarrow$ Tier 3 & +0.014 [$-$0.015, +0.043] $\mu_\mathrm{B}$/f.u. & 9/10 & none \\
\quad XGBoost & Tier 1 $\rightarrow$ Tier 2 + site count & +0.019 [$-$0.007, +0.046] $\mu_\mathrm{B}$/f.u. & 8/10 & borderline gain \\
\quad XGBoost & Tier 2 + site count $\rightarrow$ Tier 3 & +0.004 [$-$0.026, +0.033] $\mu_\mathrm{B}$/f.u. & 5/10 & none \\
\addlinespace
\multicolumn{5}{@{}l@{}}{\textit{per formula unit, m>0 only}} \\
\quad Ridge / logistic & Tier 1 $\rightarrow$ Tier 2 & +0.018 [+0.013, +0.023] $\mu_\mathrm{B}$/f.u. & 10/10 & stable gain \\
\quad Ridge / logistic & Tier 1 $\rightarrow$ Tier 3 & +0.053 [+0.036, +0.071] $\mu_\mathrm{B}$/f.u. & 10/10 & stable gain \\
\quad Ridge / logistic & Tier 2 $\rightarrow$ Tier 3 & +0.035 [+0.015, +0.054] $\mu_\mathrm{B}$/f.u. & 10/10 & stable gain \\
\quad Ridge / logistic & Tier 1 $\rightarrow$ Tier 2 + site count & +0.022 [+0.017, +0.027] $\mu_\mathrm{B}$/f.u. & 10/10 & stable gain \\
\quad Ridge / logistic & Tier 2 + site count $\rightarrow$ Tier 3 & +0.031 [+0.011, +0.051] $\mu_\mathrm{B}$/f.u. & 10/10 & stable gain \\
\quad Gradient boosting & Tier 1 $\rightarrow$ Tier 2 & +0.011 [$-$0.009, +0.032] $\mu_\mathrm{B}$/f.u. & 7/10 & borderline gain \\
\quad Gradient boosting & Tier 1 $\rightarrow$ Tier 3 & +0.031 [+0.014, +0.049] $\mu_\mathrm{B}$/f.u. & 10/10 & stable gain \\
\quad Gradient boosting & Tier 2 $\rightarrow$ Tier 3 & +0.020 [+0.003, +0.037] $\mu_\mathrm{B}$/f.u. & 10/10 & stable gain \\
\quad Gradient boosting & Tier 1 $\rightarrow$ Tier 2 + site count & +0.017 [$-$0.003, +0.038] $\mu_\mathrm{B}$/f.u. & 9/10 & borderline gain \\
\quad Gradient boosting & Tier 2 + site count $\rightarrow$ Tier 3 & +0.014 [$-$0.000, +0.029] $\mu_\mathrm{B}$/f.u. & 9/10 & none \\
\quad XGBoost & Tier 1 $\rightarrow$ Tier 2 & +0.009 [$-$0.013, +0.031] $\mu_\mathrm{B}$/f.u. & 7/10 & none \\
\quad XGBoost & Tier 1 $\rightarrow$ Tier 3 & +0.024 [+0.012, +0.037] $\mu_\mathrm{B}$/f.u. & 10/10 & stable gain \\
\quad XGBoost & Tier 2 $\rightarrow$ Tier 3 & +0.015 [$-$0.013, +0.043] $\mu_\mathrm{B}$/f.u. & 9/10 & none \\
\quad XGBoost & Tier 1 $\rightarrow$ Tier 2 + site count & +0.019 [$-$0.008, +0.046] $\mu_\mathrm{B}$/f.u. & 8/10 & borderline gain \\
\quad XGBoost & Tier 2 + site count $\rightarrow$ Tier 3 & +0.005 [$-$0.023, +0.033] $\mu_\mathrm{B}$/f.u. & 5/10 & none \\
\addlinespace
\multicolumn{5}{@{}l@{}}{\textit{per atom (moment / sites), all rows}} \\
\quad Gradient boosting & Tier 1 $\rightarrow$ Tier 3 & +2.9 [+1.5, +4.4] \% of MAE & 9/10 & stable gain \\
\quad Ridge / logistic & Tier 1 $\rightarrow$ Tier 3 & +6.0 [+5.1, +6.8] \% of MAE & 10/10 & stable gain \\
\quad XGBoost & Tier 1 $\rightarrow$ Tier 3 & +2.8 [+1.6, +3.9] \% of MAE & 10/10 & stable gain \\
\bottomrule
\end{tabularx}
\par\vspace{4pt}\raggedright\noindent Per formula unit: Nadeau-Bengio 95\% intervals, all rows and non-zero moments only. Per atom (same 10,821 rows and splits, earlier analysis with the corrected same-calculation atom count): relative MAE gain with a descriptive across-split t interval.
\end{table}


\begin{table}
\caption{Prior ML results for magnetic ordering and Curie temperature (as reported; read from the versions noted)}
\label{tab:tableS9_prior_results}
\tablefont
\begin{tabularx}{\textwidth}{@{}>{\raggedright\arraybackslash}X>{\raggedright\arraybackslash}Xl>{\raggedright\arraybackslash}X>{\raggedright\arraybackslash}X@{}}
\toprule
\textbf{Study} & \textbf{Task, data} & \textbf{Split} & \textbf{Headline result} & \textbf{Unseen-chemistry test} \\
\midrule
Nelson \& Sanvito 2019~\cite{nelson2019curie} & \Tc{}, $\sim$2,500 experimental & single partition & MAE 57 K, $R^2$ 0.87 & qualitative only \\
Long et al.\ 2021~\cite{long2021ferro} & FM/AFM + \Tc{}, 2,805 intermetallics & 90/10 & accuracy 86\%; \Tc{} MAE 58 K, $R^2$ 0.92 & no \\
Merker et al.\ 2022~\cite{merker2022classifiers} & ordering (MP labels), 34,856 & random 80/10/10 & accuracy 77.8\%, AFM F1 0.69 & no \\
Belot et al.\ 2023~\cite{belot2023curie} & \Tc{}, 2,557 + 3,194 & random 1/3 & MAE 69 K; other dataset 90 K & cross-dataset (overlapping) \\
Itani et al.\ 2025 (NEMAD)~\cite{itani2025nemad} & ordering 35,037; \Tc{} 15,577 & stratified random & accuracy 0.91, AFM F1 0.80; MAE 56 K, $R^2$ 0.87 & no \\
Orang et al.\ 2026~\cite{abediorang2026curie} & \Tc{}, 2,500 curated & random 80/20 & $R^2$ 0.87--0.91 (0.66 uncurated) & no \\
Acosta et al.\ 2022~\cite{acosta2022ordering2d} & 2D magnets, M/NM and FM/AFM & OOB; external set & accuracy 92.5\%; FM/AFM 87--90\% & external, 54 materials \\
This work & ordering 79,125; \Tc{} 29,740 & formula-, system-, provenance-grouped & macro-F1 0.73; \Tc{} MAE 76 K & yes (Extended Data Tables 5 and 6) \\
\bottomrule
\end{tabularx}
\par\vspace{4pt}\noindent {Numbers as reported by each study, on its own data and splits; they are not directly comparable with each other or with ours. Nelson \& Sanvito, Long, Belot, NEMAD, Orang and Acosta were read as arXiv versions (results/revision/research/prior\_results\_table.md lists sources and locations). Court \& Cole 2020~\cite{court2020phase} report family-specific N\'eel-temperature predictions and no dataset-wide metric.}
\end{table}

\begin{table}
\caption{Licences of the upstream data and the redistribution plan}
\label{tab:tableS10_licences}
\tablefont
\begin{tabularx}{\textwidth}{@{}>{\raggedright\arraybackslash}X>{\raggedright\arraybackslash}X>{\raggedright\arraybackslash}X>{\raggedright\arraybackslash}X@{}}
\toprule
\textbf{Source} & \textbf{Licence / terms} & \textbf{Redistribution of derived values} & \textbf{Plan} \\
\midrule
Materials Project & CC BY 4.0 (terms of use) & yes, with attribution & ship values \\
OQMD (jarvis-tools mirror) & CC BY 4.0 & yes, with attribution & ship values \\
JARVIS-DFT (Figshare) & CC BY 4.0 & yes, with attribution & ship values \\
Kelly \& Cole 2025 (Figshare) & CC BY 4.0 & yes, with attribution & ship values \\
NEMAD (Zenodo release) & CC BY 4.0 & yes, with attribution & ship values \\
AFLOW (AFLUX) & scientific, non-commercial use only & not addressed & identifiers + script, or non-commercial release \\
MAGNDATA (Bilbao server) & no public licence; conditional access & not granted publicly & identifiers + script, or permission \\
Two experimental compilations & origin undocumented & unknown & identifiers only until origin is confirmed \\
External benchmark (malaei/Magnetic\_data) & no licence & no & retrieval script with SHA-256 check \\
\bottomrule
\end{tabularx}
\par\vspace{4pt}\noindent {Evidence (quoted terms and URLs) in results/revision/research/data\_licences.md. This is a report of the terms, not legal advice; the remaining decisions rest with the authors. \todo{Implement the plan for AFLOW, MAGNDATA, the compilations and the external table before deposit.}}
\end{table}

\begin{table}
\caption{Coverage (\%) of 90\% Curie-temperature intervals in every subgroup evaluated}
\label{tab:tableS11_curie_subgroups}
\tablefont
\begin{tabularx}{\textwidth}{@{}>{\raggedright\arraybackslash}Xrrrr@{}}
\toprule
\textbf{Subgroup} & \textbf{Median formulas} & \textbf{Split} & \textbf{Normalized} & \textbf{CQR} \\
\midrule
All test materials & 2,103 & 89.3 & 90.3 & 89.4 \\
Farthest distance decile & 232 & 86.6 & 88.2 & 89.7 \\
Family: Actinide & 47 & 96.2 & 88.9 & 89.6 \\
Family: Intermetallic & 1,137 & 92.9 & 91.4 & 92.9 \\
Family: Oxide & 653 & 92.3 & 89.9 & 91.5 \\
Family: Physics-complete & 81 & 70.2 & 87.0 & 74.5 \\
Family: Rare Earth & 1,117 & 97.0 & 91.8 & 95.5 \\
Family: Transition Metal & 2,020 & 88.9 & 90.1 & 89.2 \\
Hardest sigma decile (defined by the normalized method itself) & 84 & 54.2 & 88.5 & 61.0 \\
Source: text-mined database & 544 & 78.2 & 86.0 & 79.3 \\
Source: curated compilation (anisotropy) & 461 & 92.7 & 92.4 & 92.4 \\
Source: curated compilation (magnetic materials) & 1,534 & 95.4 & 92.2 & 94.9 \\
Top decile of predicted value & 117 & 65.5 & 87.4 & 68.6 \\
\bottomrule
\end{tabularx}
\par\vspace{4pt}\raggedright\noindent Mean over 10 prespecified splits. Constant-width methods over-cover several families (rare earth, actinide, intermetallic) because their intervals are wide everywhere; normalized conformal stays within about 2 points of nominal except for text-mined labels and physics-complete materials.
\end{table}


\begin{table}
\caption{Ordering labels of Materials Project against experimental sources, and transfer between them}
\label{tab:tableS12_provenance}
\tablefont
\begin{tabularx}{\textwidth}{@{}>{\raggedright\arraybackslash}Xrrrrr@{}}
\toprule
\multicolumn{6}{@{}l@{}}{\textit{a, Labels of the 949 formulas labelled by both (rows: experimental label; columns: Materials Project label)}} \\
\midrule
\textbf{Experimental label} & \textbf{AFM} & \textbf{FM} & \textbf{FiM} & \textbf{NM} & \textbf{Total} \\
AFM & 49 & 336 & 47 & 76 & 508 \\
FM & 6 & 189 & 124 & 47 & 366 \\
FiM & 2 & 41 & 27 & 4 & 74 \\
NM & 0 & 0 & 1 & 0 & 1 \\
\midrule
\multicolumn{6}{@{}l@{}}{\textit{b, Transfer of XGBoost models on the FM, AFM and FiM classes (balanced accuracy and recall per class)}} \\
\midrule
\textbf{Training $\rightarrow$ test labels, probabilities} & \textbf{Balanced accuracy} & \textbf{FM} & \textbf{AFM} & \textbf{FiM} & \textbf{Test rows} \\
Materials Project $\rightarrow$ experimental, raw & 0.337 [0.273, 0.396] & 0.59 & 0.07 & 0.35 & 2,772 \\
Materials Project $\rightarrow$ experimental, prior-corrected & 0.314 [0.247, 0.381] & 0.67 & 0.05 & 0.22 & 2,772 \\
Experimental $\rightarrow$ Materials Project, raw & 0.384 [0.366, 0.401] & 0.19 & 0.90 & 0.06 & 4,377 \\
Experimental $\rightarrow$ Materials Project, prior-corrected & 0.400 [0.370, 0.430] & 0.60 & 0.57 & 0.04 & 4,377 \\
Experimental $\rightarrow$ experimental, raw & 0.769 [0.727, 0.811] & 0.89 & 0.90 & 0.52 & 2,772 \\
Materials Project $\rightarrow$ Materials Project, raw & 0.591 [0.562, 0.618] & 0.91 & 0.31 & 0.56 & 4,377 \\
\bottomrule
\end{tabularx}
\par\vspace{4pt}\raggedright\noindent {a, Most frequent label of each provenance per formula; agreement 28\%, Cohen's $\kappa$ 0.00 [$-$0.03, 0.03]. b, 95\% intervals from a bootstrap over formulas; prior-corrected = probabilities corrected to the class prior of the training partition of the target domain. Null from 20 refits on permuted training labels: median 0.342 and 97.5th percentile 0.374 (Materials Project $\rightarrow$ experimental), 97.5th percentile 0.351 (experimental $\rightarrow$ Materials Project). Median composition distance of the test formulas to the nearest training formula: 0.081 (Materials Project $\rightarrow$ experimental), 0.169 (experimental $\rightarrow$ Materials Project), 0.077 (within Materials Project) and 0.033 (within experimental). Gradient boosting gives the same picture (results/revision/a2\_provenance/transfer.csv).}
\end{table}

\section*{Reproducibility notes}

\paragraph{Package and checks.} The reproducibility package (paper/; README first) contains the frozen datasets, splits, results, code, tests and the table and figure builders. SHA-256 checksums cover all data and result files.

\paragraph{Later analyses.} The analyses planned in writing after the first round of checks are in the package folders \path{scripts/pipeline/revision} and \path{results/revision}, with the analysis plan and deviations log in results/revision/PLAN.md. Six planned analyses were dropped for run time before any of their results were seen:
\begin{itemize}
\item the magnetic-physics test for ordering and for the moment;
\item the Curie-temperature comparison restricted to records whose source reports the same space group as the computed polymorph;
\item the matched leave-chemical-system-out comparison;
\item seed averaging for the ordering XGBoost models;
\item a rerun of the moment intervals with per-row output (the saved subgroups of the earlier run are used);
\item a refit for the ordering decomposition (its quantities come from saved probabilities).
\end{itemize}
The categories of the title-based screen were extended during classification (Methods).

\paragraph{XGBoost thread count.} XGBoost results depend on the thread count, which acts like a random seed. Exact reproduction requires 16 threads; conclusions are stated across splits and realizations. Gradient boosting and Ridge do not depend on the thread count.

\paragraph{Runtime.} On a 16-thread, 16-GB workstation:
\begin{itemize}
\item tier fits: about 1--3 minutes per split and target;
\item uncertainty runs: 2--10 minutes per split;
\item external-benchmark analyses: minutes.
\end{itemize}
The full set of later analyses runs in about 3 hours when run sequentially.

{\singlespacing
\bibliography{references}
}
\end{document}
